\documentclass[11pt]{article}
\usepackage[utf8]{inputenc}
\usepackage[T1]{fontenc}
\usepackage[spanish]{babel}
\usepackage{amsmath, amssymb, amsthm, mathtools}
\usepackage{braket}
\usepackage{dsfont}
\usepackage[margin=2.5cm]{geometry}
\usepackage{hyperref}

\newcommand{\ii}{\mathrm{i}}
\newcommand{\ee}{\mathrm{e}}
\newcommand{\adag}{\hat a^\dagger}
\newcommand{\sigmap}{\hat\sigma_+}
\newcommand{\sigmam}{\hat\sigma_-}
\newcommand{\nnop}{\hat n}
\newcommand{\Hint}{\hat H_{\mathrm{int}}}
\newcommand{\Omn}{\tilde\Omega_n}

\title{Modelo de Jaynes-Cummings-Paul}
\author{Ricardo Del Rosal Gardu\~no \\ \small Doctorado en F\'isica, BUAP}
\date{Notas a partir de Schleich, \emph{Quantum Optics in Phase Space}}

\begin{document}
\maketitle

\tableofcontents
\bigskip

\section{Introducci\'on}

El modelo de Jaynes-Cummings-Paul describe la interacci\'on coherente entre un
\'atomo de dos niveles y un \'unico modo cuantizado del campo electromagn\'etico.
El sistema contiene el m\'inimo de ingredientes necesarios para estudiar QED de
cavidad: un grado de libertad bos\'onico, un grado de libertad at\'omico de dos
niveles y un acoplamiento dipolar que permite intercambiar una excitaci\'on entre
materia y campo.

Se considera un estado at\'omico excitado $\ket{a}$ y un estado base $\ket{b}$.
El campo se expande en la base de Fock $\{\ket{n}\}_{n=0}^{\infty}$. Por tanto,
un estado inicial factorizado se escribe como
\begin{equation}
  \ket{\psi(0)}
  =
  \ket{\psi_{\mathrm{at}}(0)}\otimes\ket{\psi_{\mathrm{c}}(0)}
  =
  \sum_{n=0}^{\infty} W_n
  \bigl(C_a\ket{a}+C_b\ket{b}\bigr)\ket{n},
\end{equation}
donde $W_n$ son las amplitudes del campo, mientras que $C_a$ y $C_b$ son las
amplitudes at\'omicas. La normalizaci\'on exige
\begin{equation}
  \sum_{n=0}^{\infty}|W_n|^2=1,
  \qquad
  |C_a|^2+|C_b|^2=1.
\end{equation}

El objetivo de la nota es obtener la evoluci\'on temporal de este estado. Se
presentan tres rutas equivalentes: exponenciaci\'on directa del operador de
evoluci\'on en resonancia, soluci\'on de las ecuaciones de Rabi fuera de
resonancia y diagonalizaci\'on matricial por bloques. Las tres rutas llevan a la
misma f\'isica: el sistema se descompone en subespacios bidimensionales
$\{\ket{a,n},\ket{b,n+1}\}$ dentro de los cuales el \'atomo y el campo
intercambian una excitaci\'on de manera reversible.

\section{Hamiltoniano de interacci\'on}

En el marco de interacci\'on, despu\'es de la aproximaci\'on dipolar y de la
aproximaci\'on de onda rotante, el hamiltoniano toma la forma
\begin{equation}
  \boxed{
  \Hint(t)
  =
  \hbar g
  \left(
  \ee^{\ii\Delta t}\adag\sigmam
  +
  \ee^{-\ii\Delta t}\sigmap \hat a
  \right).
  }
\end{equation}
Aqu\'i
\begin{equation}
  \Delta=\omega_a-\omega_c,
  \qquad
  \sigmap=\ket{a}\bra{b},
  \qquad
  \sigmam=\ket{b}\bra{a}.
\end{equation}
El operador $\adag\sigmam$ describe emisi\'on: el \'atomo baja de
$\ket{a}$ a $\ket{b}$ y el campo gana un fot\'on. El operador
$\sigmap\hat a$ describe absorci\'on: el campo pierde un fot\'on y el \'atomo
sube de $\ket{b}$ a $\ket{a}$. Los t\'erminos contra-rotantes
$\sigmap\adag$ y $\sigmam\hat a$ se descartan en la aproximaci\'on de onda
rotante porque oscilan r\'apidamente en el marco de interacci\'on y su
contribuci\'on promedio es despreciable cuando $g\ll \omega_a,\omega_c$ y el
sistema est\'a cerca de resonancia.

La ecuaci\'on de Schr\"odinger en este marco es
\begin{equation}
  \ii\hbar\frac{d}{dt}\ket{\psi(t)}
  =
  \Hint(t)\ket{\psi(t)}.
\end{equation}

\section{Caso resonante}

En resonancia exacta, $\Delta=0$, el hamiltoniano de interacci\'on deja de
depender expl\'icitamente del tiempo:
\begin{equation}
  \Hint^{(\mathrm{res})}
  =
  \hbar g
  \left(\adag\sigmam+\sigmap\hat a\right).
\end{equation}
Por tanto, la soluci\'on formal es
\begin{equation}
  \ket{\psi(t)}
  =
  \hat U(t,0)\ket{\psi(0)},
  \qquad
  \hat U(t,0)
  =
  \exp\left[
  -\ii gt
  \left(\adag\sigmam+\sigmap\hat a\right)
  \right].
\end{equation}
El problema consiste entonces en calcular la exponencial del operador
\begin{equation}
  \hat V=\adag\sigmam+\sigmap\hat a.
\end{equation}

\section{Exponenciaci\'on directa}

La exponencial se desarrolla en serie de potencias:
\begin{equation}
  \hat U(t,0)
  =
  \sum_{k=0}^{\infty}\frac{(-\ii gt)^k}{k!}\hat V^k.
\end{equation}
Para reconocer las series de seno y coseno conviene separar potencias pares e
impares:
\begin{equation}
  \hat U(t,0)
  =
  \sum_{k=0}^{\infty}\frac{(-\ii gt)^{2k}}{(2k)!}\hat V^{2k}
  +
  \sum_{k=0}^{\infty}\frac{(-\ii gt)^{2k+1}}{(2k+1)!}\hat V^{2k+1}.
\end{equation}
La estructura de todas las potencias se obtiene a partir de $\hat V^2$.

\subsection{C\'alculo de \texorpdfstring{$\hat V^2$}{V al cuadrado}}

Desarrollando el cuadrado,
\begin{align}
  \hat V^2
  &=
  \left(\adag\sigmam+\sigmap\hat a\right)
  \left(\adag\sigmam+\sigmap\hat a\right) \notag\\
  &=
  \adag\sigmam\adag\sigmam
  +
  \adag\sigmam\sigmap\hat a
  +
  \sigmap\hat a\adag\sigmam
  +
  \sigmap\hat a\sigmap\hat a .
\end{align}
Los operadores at\'omicos y de campo conmutan entre s\'i porque act\'uan en
espacios de Hilbert distintos. Adem\'as,
\begin{equation}
  \sigmam^2=0,
  \qquad
  \sigmap^2=0,
\end{equation}
ya que en un sistema de dos niveles no existe una doble bajada ni una doble
subida. Por eso el primer y el \'ultimo t\'ermino se anulan. Queda
\begin{equation}
  \hat V^2
  =
  \adag\hat a\,\sigmam\sigmap
  +
  \hat a\adag\,\sigmap\sigmam.
\end{equation}
Usando
\begin{equation}
  \sigmap\sigmam=\ket{a}\bra{a},
  \qquad
  \sigmam\sigmap=\ket{b}\bra{b},
\end{equation}
se obtiene
\begin{equation}
  \hat V^2
  =
  \hat a\adag\,\ket{a}\bra{a}
  +
  \adag\hat a\,\ket{b}\bra{b}.
\end{equation}
En la base at\'omica ordenada $\{\ket{a},\ket{b}\}$, esto se escribe como
\begin{equation}
  \hat V^2
  =
  \begin{pmatrix}
    \hat a\adag & 0\\
    0 & \adag\hat a
  \end{pmatrix}
  =
  \begin{pmatrix}
    \nnop+1 & 0\\
    0 & \nnop
  \end{pmatrix},
\end{equation}
donde $\nnop=\adag\hat a$.

\subsection{Potencias pares e impares}

La forma diagonal de $\hat V^2$ permite obtener por inducci\'on
\begin{equation}
  \hat V^{2k}
  =
  \begin{pmatrix}
    (\nnop+1)^k & 0\\
    0 & \nnop^k
  \end{pmatrix}.
\end{equation}
Para las potencias impares se multiplica por $\hat V$. Como
\begin{equation}
  \hat V
  =
  \begin{pmatrix}
    0 & \hat a\\
    \adag & 0
  \end{pmatrix},
\end{equation}
se tiene
\begin{equation}
  \hat V^{2k+1}
  =
  \begin{pmatrix}
    0 & \hat a\,\nnop^k\\
    \adag(\nnop+1)^k & 0
  \end{pmatrix}.
\end{equation}
La identidad
\begin{equation}
  \hat a\,f(\nnop)=f(\nnop+1)\hat a
\end{equation}
permite escribir el elemento superior derecho como
\begin{equation}
  \hat a\,\nnop^k=(\nnop+1)^k\hat a.
\end{equation}
Por tanto,
\begin{equation}
  \hat V^{2k+1}
  =
  \begin{pmatrix}
    0 & (\nnop+1)^k\hat a\\
    \adag(\nnop+1)^k & 0
  \end{pmatrix}.
\end{equation}

\subsection{Suma de las series}

La parte par de la exponencial es
\begin{equation}
  \sum_{k=0}^{\infty}
  \frac{(-1)^k(gt)^{2k}}{(2k)!}
  \hat V^{2k}
  =
  \begin{pmatrix}
    \cos(gt\sqrt{\nnop+1}) & 0\\
    0 & \cos(gt\sqrt{\nnop})
  \end{pmatrix}.
\end{equation}
La parte impar contiene la serie del seno. Para verla, se reescribe
$(\nnop+1)^k$ como
$(\sqrt{\nnop+1})^{2k+1}/\sqrt{\nnop+1}$. Entonces
\begin{equation}
  \sum_{k=0}^{\infty}
  \frac{(-1)^k(gt)^{2k+1}}{(2k+1)!}
  \hat V^{2k+1}
  =
  \begin{pmatrix}
    0 &
    \dfrac{\sin(gt\sqrt{\nnop+1})}{\sqrt{\nnop+1}}\hat a\\[10pt]
    \adag\dfrac{\sin(gt\sqrt{\nnop+1})}{\sqrt{\nnop+1}} &
    0
  \end{pmatrix}.
\end{equation}
Para el elemento inferior izquierdo se usa
\begin{equation}
  \adag f(\nnop+1)=f(\nnop)\adag,
\end{equation}
lo cual da
\begin{equation}
  \adag\dfrac{\sin(gt\sqrt{\nnop+1})}{\sqrt{\nnop+1}}
  =
  \dfrac{\sin(gt\sqrt{\nnop})}{\sqrt{\nnop}}\adag.
\end{equation}
La divisi\'on por $\sqrt{\nnop}$ no produce problema sobre el vac\'io porque
siempre aparece multiplicada por $\adag$ y la funci\'on se entiende por su
serie de potencias.

Finalmente,
\begin{equation}
  \boxed{
  \hat U(t,0)
  =
  \begin{pmatrix}
    \cos(gt\sqrt{\nnop+1}) &
    -\ii\,\dfrac{\sin(gt\sqrt{\nnop+1})}{\sqrt{\nnop+1}}\hat a\\[12pt]
    -\ii\,\dfrac{\sin(gt\sqrt{\nnop})}{\sqrt{\nnop}}\adag &
    \cos(gt\sqrt{\nnop})
  \end{pmatrix}.
  }
\end{equation}
La diagonal contiene amplitudes de permanencia. La antidiagonal contiene
amplitudes de transici\'on acompa\~nadas de creaci\'on o destrucci\'on de un
fot\'on. La frecuencia que aparece no es una constante \'unica, sino
$g\sqrt{n+1}$ o $g\sqrt n$, seg\'un el subespacio considerado.

\section{Acci\'on sobre estados producto}

Aplicando el operador anterior a los estados de la base producto se obtiene
\begin{equation}
  \hat U(t,0)\ket{a,n}
  =
  \cos(gt\sqrt{n+1})\ket{a,n}
  -
  \ii\sin(gt\sqrt{n+1})\ket{b,n+1},
\end{equation}
\begin{equation}
  \hat U(t,0)\ket{b,n}
  =
  \cos(gt\sqrt n)\ket{b,n}
  -
  \ii\sin(gt\sqrt n)\ket{a,n-1}.
\end{equation}
La primera expresi\'on muestra que $\ket{a,n}$ se acopla solamente con
$\ket{b,n+1}$. La segunda muestra que $\ket{b,n}$ se acopla solamente con
$\ket{a,n-1}$. Para $n=0$, el t\'ermino $\ket{a,-1}$ no existe y su coeficiente
es cero porque $\sin(0)=0$. As\'i, $\ket{b,0}$ es un estado oscuro de la
interacci\'on:
\begin{equation}
  \Hint^{(\mathrm{res})}\ket{b,0}=0.
\end{equation}

Para un estado inicial general,
\begin{align}
  \boxed{
  \begin{aligned}
  \ket{\psi(t)}
  =
  \sum_{n=0}^{\infty}W_n
  \Bigl\{
  &C_a
  \left[
  \cos(gt\sqrt{n+1})\ket{a,n}
  -
  \ii\sin(gt\sqrt{n+1})\ket{b,n+1}
  \right]\notag\\
  &+
  C_b
  \left[
  \cos(gt\sqrt n)\ket{b,n}
  -
  \ii\sin(gt\sqrt n)\ket{a,n-1}
  \right]
  \Bigr\}.
  \end{aligned}
  }
\end{align}
Si se desea agrupar por los pares $\{\ket{a,n},\ket{b,n+1}\}$, se cambia el
\'indice en los t\'erminos adecuados y se obtiene, para $n\geq 0$,
\begin{equation}
  \Psi_{a,n}(t)
  =
  W_n C_a\cos(gt\sqrt{n+1})
  -
  \ii W_{n+1}C_b\sin(gt\sqrt{n+1}),
\end{equation}
\begin{equation}
  \Psi_{b,n+1}(t)
  =
  W_{n+1}C_b\cos(gt\sqrt{n+1})
  -
  \ii W_n C_a\sin(gt\sqrt{n+1}).
\end{equation}
Estas dos amplitudes son la forma resonante que se recupera despu\'es a partir
de las ecuaciones de Rabi generales.

\section{Campo inicial en un estado de Fock}

Si el campo inicia en un estado de n\'umero definido $\ket{n_0}$, entonces
$W_n=\delta_{n,n_0}$. Para un \'atomo inicialmente excitado, $C_a=1$ y
$C_b=0$, la evoluci\'on queda
\begin{equation}
  \ket{\psi(t)}
  =
  \cos(gt\sqrt{n_0+1})\ket{a,n_0}
  -
  \ii\sin(gt\sqrt{n_0+1})\ket{b,n_0+1}.
\end{equation}
Las probabilidades at\'omicas son
\begin{equation}
  P_a(t)=\cos^2(gt\sqrt{n_0+1}),
  \qquad
  P_b(t)=\sin^2(gt\sqrt{n_0+1}).
\end{equation}
Para el vac\'io, $n_0=0$, se obtiene la oscilaci\'on de Rabi del vac\'io:
\begin{equation}
  \ket{\psi(t)}
  =
  \cos(gt)\ket{a,0}
  -
  \ii\sin(gt)\ket{b,1}.
\end{equation}
Esta es la oscilaci\'on de Rabi del vac\'io, la cual puede entenderse como
emisi\'on espont\'anea hacia un \'unico modo resonante de la cavidad, seguida de
reabsorci\'on coherente. A diferencia de la emisi\'on espont\'anea en espacio
libre, la excitaci\'on no se dispersa en un continuo de modos, sino que
permanece en el sistema \'atomo-campo y oscila entre ambos.

\section{Ecuaciones de Rabi fuera de resonancia}

Para $\Delta\neq0$, el hamiltoniano en el marco de interacci\'on depende del
tiempo. Aun as\'i, cada subespacio
\begin{equation}
  \mathcal H_n=\mathrm{span}\{\ket{a,n},\ket{b,n+1}\}
\end{equation}
permanece cerrado. La acci\'on del hamiltoniano es
\begin{equation}
  \Hint(t)\ket{a,n}
  =
  \hbar g\,\ee^{\ii\Delta t}\sqrt{n+1}\ket{b,n+1},
\end{equation}
\begin{equation}
  \Hint(t)\ket{b,n+1}
  =
  \hbar g\,\ee^{-\ii\Delta t}\sqrt{n+1}\ket{a,n}.
\end{equation}
Escribiendo
\begin{equation}
  \ket{\psi(t)}
  =
  \sum_{n=0}^{\infty}
  \left[
  \Psi_{a,n}(t)\ket{a,n}
  +
  \Psi_{b,n+1}(t)\ket{b,n+1}
  \right],
\end{equation}
y proyectando la ecuaci\'on de Schr\"odinger, se obtiene el sistema acoplado
\begin{equation}
  \dot\Psi_{a,n}(t)
  =
  -\ii g\sqrt{n+1}\,\ee^{-\ii\Delta t}\Psi_{b,n+1}(t),
\end{equation}
\begin{equation}
  \dot\Psi_{b,n+1}(t)
  =
  -\ii g\sqrt{n+1}\,\ee^{\ii\Delta t}\Psi_{a,n}(t).
\end{equation}

\subsection{Ecuaci\'on de segundo orden}

Derivando la primera ecuaci\'on,
\begin{align}
  \ddot\Psi_{a,n}
  &=
  -\ii g\sqrt{n+1}
  \left[
  -\ii\Delta\,\ee^{-\ii\Delta t}\Psi_{b,n+1}
  +
  \ee^{-\ii\Delta t}\dot\Psi_{b,n+1}
  \right].
\end{align}
El segundo t\'ermino se sustituye con la ecuaci\'on para
$\dot\Psi_{b,n+1}$:
\begin{equation}
  -\ii g\sqrt{n+1}\,\ee^{-\ii\Delta t}\dot\Psi_{b,n+1}
  =
  -g^2(n+1)\Psi_{a,n}.
\end{equation}
Para el primer t\'ermino se usa de nuevo
\begin{equation}
  -\ii g\sqrt{n+1}\,\ee^{-\ii\Delta t}\Psi_{b,n+1}
  =
  \dot\Psi_{a,n}.
\end{equation}
As\'i se llega a
\begin{equation}
  \ddot\Psi_{a,n}
  -
  \ii\Delta\,\dot\Psi_{a,n}
  +
  g^2(n+1)\Psi_{a,n}
  =
  0.
\end{equation}
La ecuaci\'on para $\Psi_{b,n+1}$ es an\'aloga:
\begin{equation}
  \ddot\Psi_{b,n+1}
  +
  \ii\Delta\,\dot\Psi_{b,n+1}
  +
  g^2(n+1)\Psi_{b,n+1}
  =
  0.
\end{equation}

\subsection{Ra\'ices caracter\'isticas}

Para $\Psi_{a,n}(t)=\ee^{\lambda t}$ se obtiene
\begin{equation}
  \lambda^2-\ii\Delta\lambda+g^2(n+1)=0.
\end{equation}
Las ra\'ices son
\begin{equation}
  \lambda_{\pm}
  =
  \frac{\ii\Delta}{2}
  \pm
  \ii\Omn,
  \qquad
  \boxed{
  \Omn=\frac{1}{2}\sqrt{\Delta^2+4g^2(n+1)}.
  }
\end{equation}
La cantidad $\Omn$ es la frecuencia de Rabi generalizada. El factor
$g\sqrt{n+1}$ mide el acoplamiento efectivo en el bloque $\mathcal H_n$.

\subsection{Soluci\'on completa con coeficientes iniciales}

La soluci\'on para la amplitud excitada se escribe como
\begin{equation}
  \Psi_{a,n}(t)
  =
  \ee^{\ii\Delta t/2}
  \left[
  A_+\ee^{\ii\Omn t}
  +
  A_-\ee^{-\ii\Omn t}
  \right].
\end{equation}
Las constantes se fijan con
\begin{equation}
  \Psi_{a,n}(0)=W_n C_a,
  \qquad
  \Psi_{b,n+1}(0)=W_{n+1}C_b,
\end{equation}
y con
\begin{equation}
  \dot\Psi_{a,n}(0)
  =
  -\ii g\sqrt{n+1}\,W_{n+1}C_b.
\end{equation}
Evaluando $\Psi_{a,n}(0)$:
\begin{equation}
  A_+ + A_- = W_n C_a.
\end{equation}
Evaluando la derivada en $t=0$:
\begin{equation}
  \ii\left(\frac{\Delta}{2}+\Omn\right)A_+
  +
  \ii\left(\frac{\Delta}{2}-\Omn\right)A_-
  =
  -\ii g\sqrt{n+1}\,W_{n+1}C_b.
\end{equation}
Resolver este sistema lineal y reordenar las exponenciales en cosenos y senos
da
\begin{align}
  \boxed{
  \Psi_{a,n}(t)
  =
  \ee^{\ii\Delta t/2}
  \left[
  \left(
  \cos(\Omn t)
  +
  \frac{\ii\Delta}{2\Omn}\sin(\Omn t)
  \right)W_n C_a
  -
  \ii\frac{g\sqrt{n+1}}{\Omn}\sin(\Omn t)\,W_{n+1}C_b
  \right].
  }
\end{align}
El mismo procedimiento para la amplitud del estado base produce
\begin{align}
  \boxed{
  \Psi_{b,n+1}(t)
  =
  \ee^{-\ii\Delta t/2}
  \left[
  -\ii\frac{g\sqrt{n+1}}{\Omn}\sin(\Omn t)\,W_n C_a
  +
  \left(
  \cos(\Omn t)
  -
  \frac{\ii\Delta}{2\Omn}\sin(\Omn t)
  \right)W_{n+1}C_b
  \right].
  }
\end{align}
Estas son las expresiones que aparecen al resolver las ecuaciones de Rabi en
las notas originales. La presencia simult\'anea de $W_n$ y $W_{n+1}$ tiene una
interpretaci\'on directa: el estado $\ket{a,n}$ recibe contribuciones desde el
componente inicial $\ket{a,n}$ y desde el componente inicial $\ket{b,n+1}$,
porque ambos viven en el mismo bloque de excitaci\'on total.

\subsection{Probabilidades y resonancia}

Si el sistema inicia en $\ket{a,n}$, entonces $W_nC_a=1$ y
$W_{n+1}C_b=0$. Las probabilidades son
\begin{equation}
  P_{a,n}(t)
  =
  1
  -
  \frac{g^2(n+1)}{\Omn^2}\sin^2(\Omn t),
\end{equation}
\begin{equation}
  P_{b,n+1}(t)
  =
  \frac{g^2(n+1)}{\Omn^2}\sin^2(\Omn t).
\end{equation}
En resonancia, $\Delta=0$, se tiene $\Omn=g\sqrt{n+1}$ y se recupera
\begin{equation}
  P_{a,n}(t)=\cos^2(g\sqrt{n+1}\,t),
  \qquad
  P_{b,n+1}(t)=\sin^2(g\sqrt{n+1}\,t).
\end{equation}
Fuera de resonancia la transferencia ya no es completa. El factor
$g^2(n+1)/\Omn^2$ reduce la amplitud m\'axima de oscilaci\'on y el desafine
introduce fases relativas $\ee^{\pm \ii\Delta t/2}$.

\section{R\'egimen altamente no resonante}

En el caso altamente no resonante,
\begin{equation}
  |\Delta|\gg 2g\sqrt{n+1},
\end{equation}
la frecuencia generalizada se expande como
\begin{align}
  \Omn
  &=
  \frac{|\Delta|}{2}
  \sqrt{1+\frac{4g^2(n+1)}{\Delta^2}}
  \notag\\
  &\approx
  \frac{|\Delta|}{2}
  \left[
  1+\frac{2g^2(n+1)}{\Delta^2}
  \right].
\end{align}
Si se toma $\Delta>0$ para fijar signos,
\begin{equation}
  \Omn\approx\frac{\Delta}{2}+\frac{g^2(n+1)}{\Delta}.
\end{equation}
El cociente que multiplica la transici\'on entre $\ket{a,n}$ y
$\ket{b,n+1}$ es
\begin{equation}
  \frac{g\sqrt{n+1}}{\Omn}\ll1.
\end{equation}
Por tanto, la transferencia real de poblaci\'on queda suprimida. Lo que queda
dominante es una fase dependiente de $n$. A segundo orden en $g/\Delta$, el
hamiltoniano efectivo puede escribirse como
\begin{equation}
  \boxed{
  \hat H_{\mathrm{eff}}
  =
  \frac{\hbar g^2}{\Delta}
  \left(
  \adag\hat a\,\sigmap\sigmam
  -
  \hat a\adag\,\sigmam\sigmap
  \right).
  }
\end{equation}
Aplicado a los estados producto,
\begin{equation}
  \hat H_{\mathrm{eff}}\ket{a,n}
  =
  \frac{\hbar g^2}{\Delta}n\ket{a,n},
\end{equation}
\begin{equation}
  \hat H_{\mathrm{eff}}\ket{b,n}
  =
  -\frac{\hbar g^2}{\Delta}(n+1)\ket{b,n}.
\end{equation}
El \'atomo y el campo ya no intercambian excitaciones, pero sus niveles se
desplazan dependiendo del n\'umero de fotones. Esta es la versi\'on cu\'antica
del desplazamiento ac-Stark. En QED de cavidad, este l\'imite permite leer el
estado del campo o del \'atomo mediante fases sin absorber necesariamente un
fot\'on.

\section{M\'etodo matricial y estados vestidos}

La tercera ruta consiste en diagonalizar el hamiltoniano completo en el marco
de Schr\"odinger:
\begin{equation}
  \hat H
  =
  \hbar\omega_c\adag\hat a
  +
  \frac{\hbar\omega_a}{2}\hat\sigma_z
  +
  \hbar g
  \left(\sigmap\hat a+\sigmam\adag\right),
\end{equation}
con
\begin{equation}
  \hat\sigma_z=\ket{a}\bra{a}-\ket{b}\bra{b}.
\end{equation}

\subsection{Conservaci\'on del n\'umero total de excitaciones}

Se define
\begin{equation}
  \hat N
  =
  \adag\hat a+\sigmap\sigmam
  =
  \adag\hat a+\ket{a}\bra{a}.
\end{equation}
La parte libre del hamiltoniano conmuta con $\hat N$ porque ambas son
diagonales en la base producto. Para la interacci\'on,
\begin{equation}
  [\sigmap\hat a,\hat N]
  =
  [\sigmap\hat a,\adag\hat a]
  +
  [\sigmap\hat a,\sigmap\sigmam].
\end{equation}
Como
\begin{equation}
  [\hat a,\adag\hat a]=\hat a,
  \qquad
  [\sigmap,\sigmap\sigmam]=-\sigmap,
\end{equation}
se obtiene
\begin{equation}
  [\sigmap\hat a,\adag\hat a]=\sigmap\hat a,
  \qquad
  [\sigmap\hat a,\sigmap\sigmam]=-\sigmap\hat a.
\end{equation}
La suma se cancela. Del mismo modo,
$[\sigmam\adag,\hat N]=0$. Por tanto,
\begin{equation}
  \boxed{[\hat H,\hat N]=0.}
\end{equation}
El n\'umero total de excitaciones es constante y el hamiltoniano se separa en
bloques independientes.

\subsection{Bloques de excitaci\'on fija}

El sector $N=0$ contiene solamente $\ket{b,0}$. Para $N=n+1\geq1$, una base
natural es
\begin{equation}
  \{\ket{a,n},\ket{b,n+1}\}.
\end{equation}
En esta base,
\begin{equation}
  \bra{a,n}\hat H\ket{a,n}
  =
  \hbar\omega_c n+\frac{\hbar\omega_a}{2},
\end{equation}
\begin{equation}
  \bra{b,n+1}\hat H\ket{b,n+1}
  =
  \hbar\omega_c(n+1)-\frac{\hbar\omega_a}{2},
\end{equation}
\begin{equation}
  \bra{a,n}\hat H\ket{b,n+1}
  =
  \bra{b,n+1}\hat H\ket{a,n}
  =
  \hbar g\sqrt{n+1}.
\end{equation}
Separando un t\'ermino proporcional a la identidad,
\begin{equation}
  \hat H^{(n)}
  =
  \hbar\omega_c\left(n+\frac{1}{2}\right)\mathds{1}
  +
  \hbar
  \begin{pmatrix}
    -\Delta_c/2 & g\sqrt{n+1}\\
    g\sqrt{n+1} & \Delta_c/2
  \end{pmatrix},
\end{equation}
donde en esta secci\'on se us\'o
\begin{equation}
  \Delta_c=\omega_c-\omega_a=-\Delta.
\end{equation}
El cambio de convenci\'on de signo solo afecta la forma de escribir los
eigenvectores, no la frecuencia $\Omn$.

\subsection{Eigenvalores}

La parte no trivial del bloque satisface
\begin{equation}
  \det
  \begin{pmatrix}
    -\hbar\Delta_c/2-\lambda & \hbar g\sqrt{n+1}\\
    \hbar g\sqrt{n+1} & \hbar\Delta_c/2-\lambda
  \end{pmatrix}
  =
  0.
\end{equation}
Entonces
\begin{equation}
  \lambda^2
  -
  \frac{\hbar^2\Delta_c^2}{4}
  -
  \hbar^2g^2(n+1)
  =
  0,
\end{equation}
y
\begin{equation}
  \lambda_\pm=\pm\hbar\Omn,
  \qquad
  \Omn=\frac{1}{2}\sqrt{\Delta_c^2+4g^2(n+1)}.
\end{equation}
Los eigenvalores completos son
\begin{equation}
  \boxed{
  E_\pm^{(n)}
  =
  \hbar\omega_c\left(n+\frac{1}{2}\right)
  \pm
  \hbar\Omn.
  }
\end{equation}

\subsection{Eigenvectores}

Los eigenvectores se escriben como superposiciones de los estados desnudos
$\ket{a,n}$ y $\ket{b,n+1}$. Se define un \'angulo de mezcla $\theta_n$ por
\begin{equation}
  \tan(2\theta_n)
  =
  \frac{2g\sqrt{n+1}}{\Delta_c}.
\end{equation}
De aqu\'i,
\begin{equation}
  \cos(2\theta_n)=\frac{\Delta_c}{2\Omn},
  \qquad
  \sin(2\theta_n)=\frac{g\sqrt{n+1}}{\Omn}.
\end{equation}
Una convenci\'on usual para los estados vestidos es
\begin{equation}
  \boxed{
  \begin{aligned}
    \ket{+,n}
    &=
    \sin\theta_n\ket{a,n}
    +
    \cos\theta_n\ket{b,n+1},
    \\
    \ket{-,n}
    &=
    \cos\theta_n\ket{a,n}
    -
    \sin\theta_n\ket{b,n+1}.
  \end{aligned}
  }
\end{equation}
Estos estados ya no son puramente at\'omicos ni puramente fot\'onicos. Son
excitaciones h\'ibridas de materia y campo. En resonancia,
$\Delta_c=0$, de modo que $\theta_n=\pi/4$ y
\begin{equation}
  \ket{\pm,n}
  =
  \frac{1}{\sqrt2}
  \left(\ket{a,n}\pm\ket{b,n+1}\right),
\end{equation}
salvo una convenci\'on de signo. La separaci\'on de energ\'ia es
\begin{equation}
  E_+^{(n)}-E_-^{(n)}
  =
  2\hbar g\sqrt{n+1},
\end{equation}
que corresponde al desdoblamiento de Rabi cu\'antico.

\section{Evoluci\'on en la base vestida}

Una vez diagonalizado el bloque, la evoluci\'on es inmediata:
\begin{equation}
  \hat U^{(n)}(t)
  =
  \ee^{-\ii\omega_c(n+1/2)t}
  \left[
  \ee^{-\ii\Omn t}\ket{+,n}\bra{+,n}
  +
  \ee^{\ii\Omn t}\ket{-,n}\bra{-,n}
  \right].
\end{equation}
Al regresar a la base producto y descartar la fase global, se obtiene
\begin{equation}
  \hat U_{\mathrm{prod}}^{(n)}(t)
  =
  \begin{pmatrix}
    \cos(\Omn t)+\ii\dfrac{\Delta}{2\Omn}\sin(\Omn t) &
    -\ii\dfrac{g\sqrt{n+1}}{\Omn}\sin(\Omn t)\\[10pt]
    -\ii\dfrac{g\sqrt{n+1}}{\Omn}\sin(\Omn t) &
    \cos(\Omn t)-\ii\dfrac{\Delta}{2\Omn}\sin(\Omn t)
  \end{pmatrix},
\end{equation}
donde se volvi\'o a la convenci\'on $\Delta=\omega_a-\omega_c$. Esta matriz
coincide con la soluci\'on obtenida por las ecuaciones de Rabi. En resonancia
se reduce a la rotaci\'on simple con frecuencia $g\sqrt{n+1}$.

\section{Conexi\'on con la tesis doctoral}

El modelo de Jaynes-Cummings-Paul es el primer bloque de construcci\'on para
la QED de cavidad usada en modelos moleculares. En el modelo de
Tavis-Cummings se reemplaza un \'atomo por muchos emisores acoplados al mismo
modo de cavidad. En el modelo de Holstein-Tavis-Cummings se a\~naden modos
vibracionales internos y acoplamiento electr\'on-fon\'on. Aun cuando el espacio
de Hilbert crece, las ideas de esta nota siguen siendo las mismas: identificar
cantidades conservadas, ordenar la base por sectores de excitaci\'on,
diagonalizar bloques cuando sea posible e interpretar las frecuencias de Rabi
como intercambio coherente entre grados de libertad.

Para el trabajo doctoral, esta nota fija la base conceptual de la interacci\'on
coherente luz-materia. El siguiente paso consiste en preguntar qu\'e queda
cuando hay varios emisores, fonones y disipaci\'on. La respuesta ya no suele ser
un operador de evoluci\'on cerrado tan simple como el JCP, pero la estructura
por bloques y la interpretaci\'on de los estados vestidos siguen siendo la gu\'ia.

\end{document}
